\subsection*{The Carrier and Operations}

\begin{definitionbox}{Definition}
\begin{definition}[Two-Element Carrier]
\label{def:two-element-carrier}
The \emph{two-element carrier} is the set
\[
\mathbb{B}_2:=\{0,1\}.
\]
\end{definition}

\begin{remark*}[Standard quantified statement]
The carrier of the two-element system is the set containing exactly the integer elements 0 and 1.
\end{remark*}

\begin{remark*}[Interpretation]
This fixes the underlying finite set before any arithmetic is assigned to it.
\end{remark*}
\end{definitionbox}

\begin{dependencies}
\begin{itemize}
  \item \hyperref[def:zero-in-z]{Zero in \(\mathbb{Z}\)}
  \item \hyperref[def:one-in-z]{One in \(\mathbb{Z}\)}
\end{itemize}
\end{dependencies}

\begin{definitionbox}{Definition}
\begin{definition}[Addition on the Two-Element System]
\label{def:addition-on-two-element-system}
Addition on \(\mathbb{B}_2\) is the binary operation \(+_2\) determined by
\[
\begin{array}{c|cc}
+_2 & 0 & 1\\ \hline
0 & 0 & 1\\
1 & 1 & 0
\end{array}
\]
\end{definition}

\begin{remark*}[Standard quantified statement]
Addition on the two-element system is the binary operation whose values are given by the displayed table.
\end{remark*}

\begin{remark*}[Interpretation]
The table is the definition of addition; in particular, 1 plus 1 equals 0 in this system.
\end{remark*}
\end{definitionbox}

\begin{dependencies}
\begin{itemize}
  \item \hyperref[def:two-element-carrier]{Two-Element Carrier}
  \item \hyperref[def:binary-arithmetic-operation]{Binary Arithmetic Operation}
\end{itemize}
\end{dependencies}

\begin{definitionbox}{Definition}
\begin{definition}[Multiplication on the Two-Element System]
\label{def:multiplication-on-two-element-system}
Multiplication on \(\mathbb{B}_2\) is the binary operation \(\cdot_2\)
determined by
\[
\begin{array}{c|cc}
\cdot_2 & 0 & 1\\ \hline
0 & 0 & 0\\
1 & 0 & 1
\end{array}
\]
\end{definition}

\begin{remark*}[Standard quantified statement]
Multiplication on the two-element system is the binary operation whose values are given by the displayed table.
\end{remark*}

\begin{remark*}[Interpretation]
The table is the definition of multiplication; 1 is multiplicative identity and 0 is absorbing.
\end{remark*}
\end{definitionbox}

\begin{dependencies}
\begin{itemize}
  \item \hyperref[def:two-element-carrier]{Two-Element Carrier}
  \item \hyperref[def:binary-arithmetic-operation]{Binary Arithmetic Operation}
\end{itemize}
\end{dependencies}

\begin{definitionbox}{Definition}
\begin{definition}[Two-Element Number System]
\label{def:two-element-number-system}
The \emph{two-element number system} is
\[
\mathbb{F}_2:=(\mathbb{B}_2,+_2,\cdot_2,0,1).
\]
\end{definition}

\begin{remark*}[Standard quantified statement]
The two-element number system is the carrier \(\mathbb{B}_2\) equipped with its addition, multiplication, zero, and one.
\end{remark*}

\begin{remark*}[Interpretation]
This packages the carrier and operations as the arithmetic structure to be verified.
\end{remark*}
\end{definitionbox}

\begin{dependencies}
\begin{itemize}
  \item \hyperref[def:two-element-carrier]{Two-Element Carrier}
  \item \hyperref[def:addition-on-two-element-system]{Addition on the Two-Element System}
  \item \hyperref[def:multiplication-on-two-element-system]{Multiplication on the Two-Element System}
\end{itemize}
\end{dependencies}

\subsection*{Verification}

\begin{theorembox}{Theorem}
\begin{theorem}[The Two-Element System Is a Field]
\label{thm:two-element-system-is-field}
The two-element number system \(\mathbb{F}_2\) is a field.
\end{theorem}

\begin{remark*}[Standard quantified statement]
The structure \(\mathbb{F}_2\) satisfies the field axioms.
\end{remark*}

\begin{remark*}[Interpretation]
The two tables make the smallest field visible by direct finite verification.
\end{remark*}
\end{theorembox}


\begin{dependencies}
\begin{itemize}
  \item \hyperref[def:two-element-number-system]{Two-Element Number System}
  \item \hyperref[def:abstract-field]{Field}
\end{itemize}
\end{dependencies}

\begin{theorembox}{Theorem}
\begin{theorem}[The Two-Element System Is Not an Ordered Field]
\label{thm:two-element-system-not-ordered-field}
There is no order on \(\mathbb{B}_2\) that makes
\(\mathbb{F}_2\) an ordered field.
\end{theorem}

\begin{remark*}[Standard quantified statement]
No order on \(\mathbb{B}_2\) makes \(\mathbb{F}_2\) satisfy the ordered-field axioms.
\end{remark*}

\begin{remark*}[Interpretation]
The equation \(1+1=0\) is compatible with field arithmetic but incompatible with ordered-field arithmetic.
\end{remark*}
\end{theorembox}


\begin{dependencies}
\begin{itemize}
  \item \hyperref[thm:two-element-system-is-field]{The Two-Element System Is a Field}
  \item \hyperref[def:abstract-ordered-field]{Ordered Field}
\end{itemize}
\end{dependencies}
